\documentclass[11pt,a4paper]{article}
\usepackage[margin=26mm]{geometry}
\usepackage{amsmath,amssymb,amsthm,booktabs,hyperref}
\newtheorem{theorem}{Theorem}
\newtheorem{lemma}[theorem]{Lemma}
\newcommand{\gen}[1]{\langle #1\rangle}
\title{Spherical hypertopes from marked kernels and antipodal flags}
\author{Human Compression project}
\date{}
\begin{document}\maketitle
\begin{abstract}
We give a common argument for the spherical classification in the fixed
manuscript \emph{An exploration of locally spherical regular hypertopes}.
Maxwell's normal-subgroup theorem is retained as a named input. Preservation
of the marked Coxeter diagram then leaves only the universal complex and,
where it exists, its antipodal quotient. A single explicit flag excludes all
even $D_n$ antipodal quotients. For $E_7$ and $E_8$, the remaining geometric
condition reduces to exact orbits of fundamental weights, avoiding an
enumeration of the full groups. Complete orbit certificates accompany the
argument.
\end{abstract}

\section{Retained inputs and the marked quotient}

We use the following standard facts about a finite Coxeter system $(W,S)$.
Its geometric representation is faithful and orthogonal; a closed fundamental
chamber meets each orbit in exactly one point. The stabilizer of a face is
its indicated standard parabolic subgroup. The spherical Coxeter complex is
a flag complex, and $W$ acts regularly on its chambers. The irreducible
finite diagrams and their group orders are the usual finite Coxeter
classification, also used in Table 1 of the source. These facts are inputs,
not new classifications established by this manuscript.

The additional normal-subgroup input is Maxwell's Theorem 0.1, quoted as
Theorem 4.1 of the fixed source: if $W$ is a finite irreducible Coxeter group
and $N\triangleleft W$, then either $N=\{1,-1\}$ or $N$ is the kernel of a
graph homomorphism to a finite Coxeter group. Here a graph homomorphism sends
each simple generator to a simple generator or to the identity, and its
images contain every simple generator of the target. The full Maxwell
paper has not been independently acquired; this is an explicitly retained
classification input through the fixed source, with its stated meaning of
``graph homomorphism.''

For a regular hypertope, the distinguished chamber automorphisms are
nonidentity involutions. Their pairwise product orders are exactly the
labels of its Coxeter diagram, including order two for an absent edge.
The resulting marked group is a quotient of the Coxeter system with that
diagram. It is a C-group, and its coset geometry is the original hypertope.
These are the acquired reconstruction facts of the L052 general manuscript.
In particular, distinct distinguished generators cannot become equal.

\begin{lemma}\label{rigidity}
Let $(W,\{s_i\}_{i\in I})$ and $(W',\{t_j\}_{j\in J})$ be Coxeter systems.
Suppose a graph homomorphism $f:W\to W'$ sends every $s_i$ to a nonidentity
element, sends distinct $s_i$ to distinct elements, and preserves the order
of every product $s_i s_k$. Then $f$ is an isomorphism of marked Coxeter
systems, after a bijection of their index sets.
\end{lemma}
\begin{proof}
Write $f(s_i)=t_{\sigma(i)}$. The assumptions on images make $\sigma$
injective; the definition of graph homomorphism makes it surjective.
The order assumption says that the two Coxeter matrices are equal under
$\sigma$. The defining presentations therefore give a homomorphism in the
opposite direction, taking $t_{\sigma(i)}$ to $s_i$. Both compositions fix
every generator, so both are identity homomorphisms.
\end{proof}

\begin{theorem}\label{two}
Subject to the stated Maxwell input, a regular hypertope whose diagram is
finite and irreducible is the universal Coxeter complex or its antipodal
quotient. The latter possibility requires $-1\in W$ and requires the quotient
coset geometry to be a hypertope.
\end{theorem}
\begin{proof}
Let $G$ be the distinguished group of the hypertope. Its marked epimorphism
$W\to G$ has kernel $N$. If Maxwell's alternative expresses $N$ as the
kernel of a graph homomorphism $f:W\to W'$, then $W'\cong W/N\cong G$.
Thus the simple images of $f$ are nontrivial, distinct, and preserve every
original product order. Lemma~\ref{rigidity} gives $N=1$. The remaining
alternative is precisely $N=\{1,-1\}$. Reconstruction identifies the actual
hypertope with the coset geometry of the indicated quotient and marking.
No assumption that its proper residues are universal has been used.
\end{proof}

This argument avoids a separate inspection of the normal-subgroup tables
for each finite type. It also avoids using injectivity on a parabolic before
knowing that the corresponding residue is universal.

\section{Parabolics in the antipodal quotient}

Let $\alpha_1,\ldots,\alpha_n$ be simple roots in the faithful real
representation, and choose fundamental vectors $\omega_i$ with
$\langle\omega_i,\alpha_j\rangle=\delta_{ij}$. Their Gram matrix is the
inverse of the simple-root Gram matrix. For an irreducible finite system,
\begin{equation}\label{positive}
 c_{ij}:=\langle\omega_i,\omega_j\rangle>0\qquad(1\leq i,j\leq n).
\end{equation}
For completeness, write the positive definite simple-root Gram matrix as
$D^{1/2}(1-A)D^{1/2}$, where $D$ is positive diagonal and $A$ has nonnegative
entries and connected support. Positive definiteness implies that the
largest eigenvalue of $A$ is less than one. For a symmetric nonnegative
matrix every eigenvalue has absolute value at most its largest eigenvalue;
this follows from $|x^TAx|\leq |x|^TA|x|$. Hence
$(1-A)^{-1}=\sum_{m\geq0}A^m$. Connectivity supplies a positive entry in
some power for every pair of indices, proving \eqref{positive}.

Write $W_i=\gen{s_j:j\ne i}$. This subgroup fixes the nonzero vector
$\omega_i$, so it cannot contain $-1$. Moreover,
\begin{equation}\label{products}
 -1\notin W_iW_j\qquad(i\ne j).
\end{equation}
Indeed, if $-1=ab$ with $a\in W_i$ and $b\in W_j$, then
$a\omega_j=-\omega_j$ and $a\omega_i=\omega_i$. Orthogonality would give
$c_{ij}=-c_{ij}$, contrary to \eqref{positive}.

For rank at least three, the quotient by $\{1,-1\}$ is therefore injective
on every proper parabolic and satisfies the exact product-kernel criterion
proved in \emph{Kernels and intersections in quotients of hypertopes}.
It is a C-group on the induced generators. This conclusion does not assert
flag transitivity: the latter is precisely what fails in even type $D_n$.
For the string diagrams, the standard string C-group reconstruction theorem
does supply a regular polytope and hence a regular hypertope.

\section{Incidence and signs}

The vertices of type $i$ in the spherical Coxeter complex can be represented
by $W\omega_i$. Two vertices of distinct types are incident exactly when
they belong to one chamber. We need the following precise numerical test.

\begin{lemma}\label{dot}
For $v\in W\omega_j$, one has
$\langle\omega_i,v\rangle\leq c_{ij}$. Equality holds if and only if
$v\in W_i\omega_j$, equivalently if $v$ is incident to $\omega_i$.
\end{lemma}
\begin{proof}
If $v$ is not in the closed fundamental chamber, some simple root satisfies
$\langle v,\alpha_l\rangle<0$. Reflect in that root. This increases the
inner product with $\eta=\sum_k\omega_k$ strictly, and does not decrease
the inner product with $\omega_i$: the increase in the latter is zero
unless $l=i$, in which case it is strictly positive. The orbit is finite,
so repeating reaches its unique dominant representative $\omega_j$.
This proves the inequality. If equality held initially, no step with
$l=i$ could have occurred. All steps therefore belong to $W_i$, and
$v\in W_i\omega_j$. Conversely $W_i$ fixes $\omega_i$, so every member
of this orbit gives equality. The parabolic description of the Coxeter
complex identifies $W_i\omega_j$ as precisely the vertices of type $j$
incident to $\omega_i$.
\end{proof}

Suppose $-1\in W$. An antipodal vertex is $[v]=\{v,-v\}$. By
Lemma~\ref{dot}, two such vertices of distinct types are incident exactly
when their chosen representatives have inner product $c_{ij}$ or
$-c_{ij}$. Positivity makes the sign unambiguous.

\begin{lemma}\label{triangle}
Every flag in the antipodal quotient lifts to a flag in the spherical
complex if and only if every incident triple has positive product of
its three incidence signs. It suffices to test that there are no distinct
$i,j,k$ and $v\in W\omega_k$ such that
\begin{equation}\label{bad}
 \langle\omega_i,v\rangle=c_{ik},\qquad
 \langle\omega_j,v\rangle=-c_{jk}.
\end{equation}
\end{lemma}
\begin{proof}
Changing the sign of one representative changes the signs of its incident
pairs, and leaves the product around a triangle unchanged. If all triangle
products in a flag are positive, choose a first vertex and orient every
other representative to have positive inner product with it. The triangle
condition makes every remaining pair positive as well. These representatives
are pairwise incident in the original complex; its flag property makes them
a flag and hence a subflag of a chamber. Conversely any lifted flag admits
all-positive representatives.

If some triangle has negative product, orient its first two representatives
to make their pair positive. They are incident in the original complex,
so an element of $W$ takes them to $\omega_i,\omega_j$. Choose the sign of
the third representative to make its pairing with $\omega_i$ positive.
Its pairing with $\omega_j$ must then be negative, which is exactly
\eqref{bad}. This procedure is reversible.
\end{proof}

When the condition holds, every quotient flag lies in the image of a
chamber, and $W/\{1,-1\}$ is flag transitive. The C-group reconstruction
with flag transitivity now gives the hypertope. This is stronger than merely
checking that every face has the correct number of adjacent chambers.

\section{A flag excluding every even \texorpdfstring{$D_n$}{Dn} quotient}

Use the standard realization of $D_n$ by signed permutations with an even
number of sign changes, where $n\geq4$ is even. Its fundamental vertices
include
\[
 e_1,\qquad
 w=\tfrac12(e_1+\cdots+e_{n-1}-e_n),\qquad
 v=\tfrac12(e_1+\cdots+e_n).
\]
The central antipode belongs to this group exactly when $n$ is even.
Put $x=e_n$. The vertices $v,w$ are incident, as they are in the displayed
fundamental chamber. The vertices $x,v$ are incident because a coordinate
permutation takes $e_1$ to $x$ and fixes $v$. Likewise $-x,w$ are incident:
apply the permutation interchanging coordinates $1,n$, followed by sign
changes in those two coordinates. This group element takes $e_1$ to $-e_n$
and fixes $w$.

Consequently $[v],[w],[x]$ form a flag in the antipodal incidence system.
Its pairings have signs $+,+,-$, since
\[
 \langle v,w\rangle=(n-2)/4>0,\quad
 \langle v,x\rangle=1/2,\quad
 \langle w,x\rangle=-1/2.
\]
It cannot lift to a chamber. This precludes regularity with the proposed
marked group: if the quotient were such a regular hypertope,
its chambers would be the images of the original chambers by coset
reconstruction and chamber transitivity. The negative triple contradicts
that requirement. Thus no regular hypertope of type $D_n$ can have the
antipodal quotient as its distinguished group. This one flag handles every
even rank, without a separate chamber lower bound or induction on rank.

\section{The two exceptional incidence calculations}

For $E_7$ and $E_8$, choose the simply laced Cartan matrix $C$ with diagonal
entries two and entries minus one along the tree. Number the trivalent
node by $0$. Its three arms have vertices $(1)$, $(2,3)$, and $(4,5,6)$
for $E_7$, with $(4,5,6,7)$ as the third arm for $E_8$. All other
entries are zero. Put $d=\det C$, so $d=2$ or $1$, respectively, and
$B=dC^{-1}$. These are integral matrices with strictly positive entries.

Represent $d\omega_k$ in the simple-root basis by column $k$ of $B$.
A simple reflection changes an integer vector $z$ by
\begin{equation}\label{reflection}
 z\longmapsto z-\left(\sum_l C_{il}z_l\right)e_i.
\end{equation}
Moreover $\langle\omega_i,z\rangle=z_i$. Therefore the excluded pattern
\eqref{bad} becomes the exact pair of integer equalities
\[
 z_i=B_{ik},\qquad z_j=-B_{jk},\qquad i,j,k\text{ distinct}.
\]
Starting at column $k$, close a set under all the reflections
\eqref{reflection}. Every stored vector belongs to the orbit, by induction
on its construction. When no new vector is found, the set is invariant
under every generator and contains the starting vector, so it is the full
orbit. Thus no group-order estimate or randomized sampling enters the
coverage argument.

The complete orbit certificates have the following sizes and contain no
excluded pattern. A separate reader checks that each set contains its
initial weight and is closed under every simple reflection. These two
conditions already suffice for coverage of the required orbit.
\[
\begin{array}{c|rrrrrrrr}
 &0&1&2&3&4&5&6&7\\\hline
 E_7&10080&576&2016&126&4032&756&56&-\\
 E_8&483840&17280&69120&2160&241920&60480&6720&240
\end{array}
\]
These fifteen orbits contain $899402$ vectors in total. The reader derives
$B$ again by rational elimination from the displayed Dynkin trees. It checks
strict ordering, uniqueness, the initial weight, every reflection image,
and the absence of the excluded coordinate pattern. It does not import
the orbit construction. The resulting finite sets are supplied in full,
with deterministic compressed copies. The calculation enumerates the
fundamental-weight orbits, rather than the groups of orders $2903040$
and $696729600$.

All arithmetic is integral. Indeed $z^TCz=dB_{kk}$ for an orbit vector;
Cauchy--Schwarz gives $z_i^2\leq B_{ii}B_{kk}\leq900$.
Every stored coordinate therefore has absolute value at most thirty.
The producer checks this bound after each reflection; its intermediate
integer expressions are bounded by the displayed matrices. A deadline or
storage failure is an incomplete calculation and cannot return a complete
certificate. The reader also rejects a set with a deleted vector, a changed
coordinate, or an incorrect orbit header; the producer rejects a changed
Cartan matrix. These controls exercise coverage and exact input identity.

\section{Classification and proper residues}

For clarity, the central antipode can be decided independently of the
classification under discussion. In type $A_n$ with $n\geq2$, a coordinate
permutation on the sum-zero space cannot negate every difference $e_i-e_j$:
three distinct coordinates already give contradictory required images.
In $B_n$ every signed permutation occurs, whereas in $D_n$ only an even
number of sign changes is allowed. Thus $-1$ occurs in $B_n$ and exactly
the even $D_n$.

For $E_6$, the Cartan determinant is three. The weight lattice modulo the
root lattice consequently has order three. Each simple reflection acts
trivially on this quotient, since its difference from the identity is an
integral multiple of a simple root. Negation acts nontrivially on a group
of order three, so it cannot belong to the Weyl group. The determinant
and the usual weight/root-lattice identification are retained elementary
Coxeter facts here.

For the remaining types, the accompanying exact matrix calculation verifies
that the product of simple reflections, raised to the indicated power,
is $-I$:
\[
\begin{array}{c|rrrrr}
\text{type}&E_7&E_8&F_4&H_3&H_4\\\hline
\text{power}&9&15&6&5&15.
\end{array}
\]
The $E$ matrices are those already displayed. For $F_4$ use the chain
Cartan matrix with diagonal two, neighbouring entries minus one except
$C_{12}=-2$ (indices start at zero). For $H_3,H_4$ use the symmetric
matrix with diagonal two, first bond $-\phi$, and other chain bonds
minus one, where $\phi^2=\phi+1$. Reflection $i$ subtracts row $i$
of this matrix from row $i$ of the identity. Multiplication is exact in
$\mathbb Z[\phi]$, or in $\mathbb Z$ in the integral cases. Faithfulness
of these retained geometric representations identifies the computed
matrix with the antipode in the Coxeter group. In rank one the antipodal quotient
loses its generator. In rank two, quotienting $I_2(p)$ by a nontrivial
central antipode reduces the order of the product of its generators,
and hence changes the prescribed diagram.

The complete finite checks, Theorem~\ref{two}, the string reconstruction,
the even $D_n$ flag, and Lemma~\ref{triangle} give the following classification.
The column $a$ records whether the antipodal quotient is also admitted.
The group orders in the universal column are retained from the finite
Coxeter classification.
\[
\begin{array}{c|c|r|c}
\text{diagram}&\text{restriction}&|W|&a\\\hline
A_n&n\geq1&(n+1)!&0\\
D_n&n\geq4&2^{n-1}n!&0\\
B_n&n\geq3&2^nn!&1\\
I_2(p)&p\geq3&2p&0\\
E_6&&51840&0\\
E_7&&2903040&1\\
E_8&&696729600&1\\
F_4&&1152&1\\
H_3&&120&1\\
H_4&&14400&1
\end{array}
\]
There are exactly $1+a$ hypertopes in each row. When $a=1$, the second
object has distinguished group of order $|W|/2$. The two objects are
distinct because they have different numbers of chambers.

Every proper residue of an admitted antipodal quotient is universal
spherical. Indeed a maximal parabolic embeds, as proved above. Flag
transitivity and the C-group intersection identities identify the residue
of its corresponding face with the coset geometry of that embedded
parabolic on its own maximal parabolics. This is its original spherical
Coxeter complex. Residues of higher codimension are residues within these.
Thus proper spherical residues are a consequence of the classification,
not an assumption silently used to force the quotient kernel.

\paragraph{Attribution and limits.}
The classification conclusions and the standard Coxeter structure belong
to the earlier literature. Maxwell's normal-subgroup theorem is a substantial
retained input. The work here is the common marked rigidity argument, the
explicit antipodal flag criterion, the uniform $D_n$ obstruction, and the
exceptional orbit proof with its exact finite obligations. No historical
priority or open-conjecture solution is claimed. The Euclidean and hyperbolic
classifications in the original manuscript lie outside the scope of this note.

\begin{thebibliography}{9}
\bibitem{source} M. E. Fernandes, D. Leemans, and A. Ivi\'c Weiss,
\emph{An exploration of locally spherical regular hypertopes}.
Fixed sixteen-page revised manuscript, catalogue L033; Sections 2--4.
\bibitem{maxwell} G. Maxwell, \emph{The normal subgroups of finite and affine
Coxeter groups}, Proc. London Math. Soc. 76 (1998), 359--382.
Theorem 0.1 retained through the fixed source; primary full text not acquired.
\bibitem{components} Human Compression project,
\emph{Residue stabilizers and distinguished generators};
\emph{Kernels and intersections in quotients of hypertopes}.
Complete reviewed components acquired with their proofs and stated inputs.
\end{thebibliography}
\end{document}
